% !TEX root = A0.MiTFG.tex

\chapter{Auditoría de seguridad del ecosistema Widevine}

Este capítulo presenta la auditoría de la Fase~0 del laboratorio descrito en el capítulo anterior. Esta fase constituye la línea base del trabajo: dispone de reproducción Widevine funcional, un cliente OTT, una CDN local, un servidor de origen y un proxy de licencias, pero conserva deliberadamente rutas y decisiones de autorización insuficientes. El objetivo de la auditoría no es romper la criptografía de Widevine ni extraer claves del CDM, sino comprobar si un tercero puede consumir el contenido o utilizar la infraestructura fuera del contexto autorizado por la plataforma.

La distinción es esencial para interpretar los resultados. Widevine protege el contenido cifrado y procesa licencias dentro de un CDM compatible, mientras que la autenticación del usuario, la autorización por activo y el control de las sesiones pertenecen a la aplicación que integra el DRM \cite{w3ceme2026,widevine2024}. Por tanto, los hallazgos descritos a continuación afectan principalmente al ecosistema que rodea a Widevine. La auditoría evalúa si la plataforma mantiene una relación verificable entre identidad, activo, sesión, licencia y distribución, y no si el algoritmo de cifrado del DRM puede ser vulnerado.

\section{Metodología de auditoría}

La auditoría se realizó mediante un enfoque mixto de caja negra y caja gris. La perspectiva de caja negra permitió observar el sistema como lo haría un usuario o un tercero con acceso a las URLs públicas. La perspectiva de caja gris añadió el conocimiento de la topología Docker, las rutas de Varnish y la lógica del servidor para explicar por qué cada comportamiento era posible. La combinación de ambas evita limitar el análisis a una mera revisión de código y, al mismo tiempo, permite relacionar cada resultado observable con una decisión concreta de arquitectura.

El alcance quedó restringido a los componentes de Fase~0: cliente oficial, reproductor externo, laboratorio DRM, demostración de clave conocida, Varnish, Nginx y proxy de licencias. Se utilizaron exclusivamente usuarios, claves y activos de laboratorio. El contenido Widevine y el servicio de licencias aguas arriba corresponden a recursos públicos de prueba del ecosistema Shaka; no se emplearon contenidos comerciales, credenciales de terceros ni técnicas de extracción de claves.

La matriz experimental comenzó estableciendo el comportamiento legítimo de referencia. Se comprobó la reproducción de los dos activos con el usuario autorizado y el rechazo del usuario sin permisos. Después se repitieron las solicitudes desde clientes ajenos a la plataforma y mediante rutas alternativas. Este orden permite atribuir los resultados a las decisiones de seguridad de Fase~0 y no a una falta de funcionalidad del laboratorio.

\subsection{Caja negra}

El análisis de caja negra partió de las interfaces expuestas en el equipo anfitrión. El cliente oficial se publicó en el puerto \texttt{3000}; el reproductor externo, el laboratorio DRM y la demostración de clave conocida se publicaron respectivamente en los puertos \texttt{3001}, \texttt{3002} y \texttt{3003}. Varnish actuó como CDN en el puerto \texttt{8080}, mientras que el origen y el servicio de licencias también quedaron accesibles directamente en \texttt{8081} y \texttt{8082}.

Desde esta perspectiva se plantearon tres perfiles. El primero fue un usuario autenticado y autorizado, utilizado para establecer el comportamiento legítimo. El segundo fue un usuario autenticado sin \textit{entitlement}, con el que se comprobó si la plataforma impedía obtener la configuración del activo. El tercero fue un cliente externo sin sesión en la plataforma, que únicamente disponía de la URL del MPD y de una dirección de licencia observables.

Las pruebas verificaron códigos HTTP, estados visibles del reproductor, disponibilidad de botones, avance temporal del vídeo y posibilidad de acceder a rutas alternativas. Se consideró que una reproducción había comenzado realmente cuando el elemento de vídeo no estaba pausado, su tiempo de reproducción avanzaba y el navegador informaba de datos multimedia disponibles. Esta comprobación impidió confundir la mera aparición de un mensaje de éxito en la interfaz con una reproducción efectiva.

\subsection{Caja gris}

El análisis de caja gris revisó la configuración de contenedores, las reglas de Varnish y la lógica del servidor de Fase~0. En particular, se examinó qué servicios se publicaban en el anfitrión, qué rutas se enviaban al origen o al proxy de licencias, cómo se generaban los tokens de acceso y qué validaciones se ejecutaban antes de reenviar un \textit{challenge} Widevine.

Esta revisión mostró que la plataforma mantiene dos caminos de licencia. El camino oficial, \texttt{/platform/license}, exige un token de acceso y comprueba que el usuario disponga del activo indicado. El camino alternativo, \texttt{/license/no\_auth}, reenvía el mismo tipo de mensaje binario sin exigir identidad ni autorización. La coexistencia de ambos caminos es intencionada en Fase~0 y permite demostrar que un control correcto en la aplicación pierde eficacia si permanece disponible una ruta equivalente que no aplica la misma política.

La caja gris también permitió comprobar que los tokens de acceso son identificadores UUID almacenados en memoria. No contienen expiración, ámbito, activo, dispositivo o identificador de sesión, y permanecen aceptados hasta que el proceso se reinicia. Asimismo, el origen y el servidor de licencias se publican directamente en el anfitrión, de modo que el paso por Varnish no es obligatorio.

\subsection{Análisis de tráfico}

El tráfico se observó en tres niveles complementarios. En el navegador se registraron las peticiones de manifiesto y licencia generadas por Shaka Player. En la capa HTTP se inspeccionaron códigos de estado y cabeceras de respuesta. Finalmente, los registros de los contenedores permitieron confirmar las solicitudes recibidas por Nginx durante la descarga del MPD, la inicialización y los segmentos locales.

El flujo externo Widevine produjo una solicitud al MPD público, dos solicitudes al endpoint \texttt{/license/no\_auth} y dos respuestas binarias de licencia de 714 y 1229 bytes. El laboratorio registró después el evento \texttt{PLAYBACK\_STARTED}, y el elemento de vídeo alcanzó estado de reproducción con datos suficientes. En el activo CENC local, los logs del origen mostraron solicitudes al fichero de inicialización y a los segmentos \texttt{m4s}, sin que se produjera ninguna petición a un servidor de licencias.

Las cabeceras confirmaron además que los objetos locales atravesaban la caché de Varnish: una segunda solicitud al MPD devolvió \texttt{X-Cache: HIT}. El mismo MPD pudo obtenerse directamente desde el origen publicado en \texttt{8081}. Por tanto, el análisis no solo observó que el contenido estaba disponible, sino también que existían dos caminos de distribución: uno a través de la CDN y otro directo hacia el origen.

No se intentó interpretar el contenido criptográfico de las respuestas Widevine ni recuperar claves desde el tráfico. EME no expone las claves del CDM al código JavaScript de la página, y una captura de red no debe confundirse con una vía de exportación del material protegido \cite{w3ceme2026,patat2022}. Para los resultados de este capítulo fueron suficientes la secuencia de peticiones, los tamaños de respuesta, el estado del reproductor y los logs del servidor.

\subsection{Reproducción de ataques}

Los escenarios se reprodujeron de forma controlada y repetible. Cada prueba comenzó con los siete contenedores de Fase~0 activos y sin reinicios. Se verificó primero el comportamiento legítimo, después el rechazo del usuario sin permisos y, por último, los caminos alternativos disponibles para un cliente externo.

La reproducción se limitó a acciones que un usuario podría realizar sobre sus propias comunicaciones: observar URLs, introducirlas en otro reproductor, reutilizar un token emitido por el laboratorio y acceder a puertos publicados localmente. La demostración de clave conocida empleó un KID y una clave generados expresamente para el TFM. No se modificó el CDM, no se instrumentó Widevine L3 y no se recuperaron claves de terceros. Este límite mantiene la prueba alineada con el objetivo del trabajo: auditar la integración OTT y no desarrollar una herramienta de elusión del DRM.

\section{Análisis del flujo de licencias}

En el flujo oficial, el usuario inicia sesión mediante \texttt{/auth/login}. El servidor comprueba las credenciales de laboratorio, genera un token aleatorio y devuelve un catálogo en el que cada activo aparece como permitido o denegado. Cuando el usuario autorizado selecciona el activo Widevine, el cliente solicita su configuración mediante \texttt{/playback/config}. Shaka Player descarga después el MPD, crea una sesión EME y envía el \textit{challenge} a \texttt{/platform/license} con el token de acceso y el identificador de activo.

El servidor vuelve a comprobar la identidad y el \textit{entitlement} antes de reenviar el mensaje al servicio Widevine de pruebas. Este camino produjo reproducción real: el elemento de vídeo se mantuvo en ejecución, alcanzó un estado de datos completos y avanzó temporalmente. Por tanto, Fase~0 dispone de un flujo protegido funcional dentro de la interfaz oficial.

No obstante, el mismo proxy publica \texttt{/license/no\_auth}. Este endpoint acepta el cuerpo binario generado por el CDM y lo reenvía sin token, usuario o activo. El reproductor externo configuró \texttt{com.widevine.alpha} con esa URL, obtuvo respuestas de licencia válidas y reprodujo el mismo activo sin iniciar sesión. La licencia siguió siendo procesada correctamente por Widevine; el fallo de seguridad consistió en que la aplicación proporcionó acceso al servicio de licencia fuera de su política de negocio.

\subsection{Obtención de tokens}

El token de Fase~0 se obtiene tras enviar correo y contraseña a \texttt{/auth/login}. El valor emitido es un UUID aleatorio y opaco asociado en memoria a la dirección de correo. Este mecanismo evita que un usuario no autenticado invoque directamente la ruta oficial de plataforma, pero aporta garantías limitadas frente a robo o reutilización.

El token no incluye fecha de expiración, identificador único verificable fuera del mapa de memoria, audiencia, ámbito ni contexto de reproducción. Tampoco se rota al iniciar una reproducción o al solicitar una licencia. A diferencia de un JWT con \textit{claims} explícitos y tiempo de expiración, el servidor depende exclusivamente de que el UUID continúe presente en su estructura interna \cite{rfc7519}. Reiniciar el proceso invalida todos los tokens, pero mientras el proceso siga activo no existe caducidad automática.

Durante la revisión visual se observó además que la interfaz oficial imprime la respuesta completa de login en el bloque de estado, incluido el token de acceso. Aunque un usuario puede observar su propio token en las herramientas de red, mostrarlo directamente en el DOM incrementa la probabilidad de exposición mediante capturas de pantalla, soporte remoto o registros copiados. El token no se reproduce en esta memoria para evitar convertir la evidencia en una credencial reutilizable.

\subsection{Validaciones existentes}

Fase~0 no carece por completo de controles. La ruta oficial valida credenciales, diferencia dos cuentas y comprueba el \textit{entitlement} antes de devolver la configuración. El usuario autorizado pudo reproducir tanto el activo Widevine como el activo CENC local. El usuario sin permisos pudo autenticarse, pero recibió ambos activos como denegados y los botones de reproducción permanecieron deshabilitados.

El endpoint \texttt{/platform/license} repite la comprobación de autorización en cada \textit{challenge}. Esta segunda validación es positiva porque evita confiar únicamente en la decisión adoptada al mostrar el catálogo. También se impide solicitar una licencia Widevine para un activo declarado como ClearKey.

Sin embargo, estas validaciones se aplican solo al camino oficial. La ruta \texttt{/license/no\_auth} no exige token, y las rutas de contenido local tampoco requieren autorización. El identificador de activo se recibe en una cabecera controlada por el cliente y no se contrasta con el KID o el PSSH contenido en el mensaje Widevine. Además, no se valida dispositivo, IP, agente de usuario, sesión de reproducción, número de accesos simultáneos ni ritmo de peticiones. En consecuencia, la autorización existente es puntual y no representa una política extremo a extremo.

\subsection{Dependencias de confianza}

La prueba depende de varios componentes externos e internos. El navegador confía en Shaka Player para coordinar EME y en el CDM instalado para procesar la licencia. El MPD Widevine se descarga desde el almacenamiento público de demostración de Shaka, y el proxy local confía en el servicio CWIP para emitir la respuesta. La CDN local confía en el origen Nginx para entregar los objetos, mientras que el cliente confía en que las URLs proporcionadas por el backend corresponden al activo autorizado.

Estas dependencias delimitan lo que demuestra el experimento. El éxito del escenario externo acredita que el proxy de licencias local acepta solicitudes sin identidad y que Widevine reproduce cuando recibe una licencia válida. No acredita que se haya comprometido el servicio CWIP, el CDM o la criptografía de Widevine. De igual forma, el contenido Widevine público no atraviesa el origen local, por lo que la prueba de acceso directo a \texttt{8081} se realiza sobre activos locales y no sobre los segmentos Widevine alojados por Google.

La configuración usa además HTTP y permite CORS desde cualquier origen. CORS no constituye un mecanismo de autenticación y solo limita determinados comportamientos de navegadores; un cliente nativo o una herramienta de línea de comandos no queda sometido a esa política. Por ello, restringir CORS sería una medida de reducción de superficie, pero no resolvería por sí sola la ausencia de autorización en la licencia o el contenido.

\section{Evaluación de vulnerabilidades}

La Tabla~\ref{tab:hallazgos-fase0} resume los principales hallazgos. La severidad se asignó cualitativamente atendiendo al impacto sobre la confidencialidad del contenido, la facilidad de explotación y la necesidad de disponer de credenciales previas.

\begin{table}[H]
\centering
\caption{Resumen de hallazgos de seguridad en Fase~0}
\label{tab:hallazgos-fase0}
\begin{tabularx}{\textwidth}{p{0.10\textwidth}p{0.24\textwidth}Yp{0.13\textwidth}}
\toprule
\textbf{ID} & \textbf{Hallazgo} & \textbf{Impacto principal} & \textbf{Severidad} \\
\midrule
F0-01 & Licencia pública sin autenticación & Reproducción Widevine fuera de la plataforma sin identidad ni permiso. & Crítica \\
F0-02 & Exposición directa de origen y licencias & Elusión del frontal CDN y ampliación de la superficie accesible. & Alta \\
F0-03 & Token sin expiración ni contexto & Reutilización durante toda la vida del proceso ante una filtración. & Alta \\
F0-04 & Ausencia de control de concurrencia & Uso simultáneo y compartición de cuenta sin detección. & Alta \\
F0-05 & Contenido local sin autorización & Descarga de MPD y segmentos con solo conocer sus rutas. & Alta \\
F0-06 & CORS y transporte permisivos & Facilita integración desde terceros y expone datos de control en HTTP. & Media \\
F0-07 & Token mostrado en la interfaz & Aumenta el riesgo de divulgación accidental de la credencial. & Media \\
\bottomrule
\end{tabularx}
\end{table}

\subsection{CDN \textit{leeching}}

La CDN permite recuperar el MPD y los segmentos locales sin token. El conocimiento de una ruta como \texttt{/content/dash-known-key/stream.mpd} resulta suficiente para que cualquier cliente alcance el activo. Varnish cachea el objeto y lo sirve con \texttt{Access-Control-Allow-Origin: *}, por lo que la infraestructura puede consumir ancho de banda aunque la solicitud no proceda de la aplicación oficial.

El origen también está publicado en \texttt{8081}. El mismo MPD obtenido mediante la CDN pudo descargarse directamente desde Nginx, evitando las reglas de Varnish. Este resultado reproduce una condición frecuente de \textit{bypass}: proteger únicamente el frontal no es efectivo si permanece disponible una ruta alternativa hacia el origen. Los mecanismos de URLs o cookies firmadas descritos por proveedores de CDN persiguen precisamente que conocer una ruta no sea suficiente para acceder al objeto \cite{googlemediacdn2026,awscloudfront2026}.

Debe matizarse el alcance del resultado. El activo Widevine de Sintel está alojado por Google y sus segmentos no atraviesan la CDN local. Por tanto, el laboratorio demuestra acceso no autorizado a objetos CENC locales, exposición directa del origen y abuso del proxy de licencias, pero no una cadena completa de \textit{leeching} de segmentos Widevine servidos por la propia CDN. La evaluación de tokens de acceso a CDN sigue siendo relevante para el diseño defensivo, pero esta limitación impide presentar el escenario como una reproducción completa de un proveedor comercial \cite{simon2024cdn,rodrigues2024cdn}.

\subsection{Reutilización de sesiones}

Fase~0 no crea una sesión de reproducción diferenciada. El token de login es la única referencia de estado y puede reutilizarse para pedir configuraciones y licencias oficiales mientras permanezca en memoria. No existe \textit{logout}, lista de revocación, expiración temporal, nonce de un solo uso ni asociación entre token y una reproducción concreta.

Si el token se filtra, otro cliente puede presentar el mismo valor y actuar como el usuario hasta que el proceso se reinicie. El servidor tampoco registra un cambio de IP o dispositivo que permita detectar el traslado de la credencial. Este comportamiento no rompe la firma de un token ni el protocolo Widevine; explota la ausencia de ciclo de vida y contexto alrededor de la credencial. Los mecanismos propuestos para las fases endurecidas deberán separar el token de autenticación del token de reproducción y reducir la ventana de reutilización.

\subsection{\textit{Credential sharing}}

La compartición de credenciales es viable porque la cuenta no se vincula a dispositivos y el servidor no conserva información sobre ubicaciones o patrones de acceso. Varias personas pueden utilizar el mismo correo y contraseña y obtener tokens distintos sin que el sistema relacione esas autenticaciones entre sí. El DRM continuará emitiendo licencias válidas porque cada petición llega desde un navegador compatible y el plano de negocio no detecta que la identidad se está compartiendo.

En el laboratorio no se simularon desplazamientos geográficos reales ni poblaciones masivas de clientes. No obstante, la revisión de caja gris confirma que no existe ninguna función capaz de distinguir esos casos. Por ello, el hallazgo se fundamenta en la ausencia verificable del control y no en una inferencia sobre la criptografía. La compartición de cuentas pertenece al plano de política de servicio y requiere señales adicionales de cuenta, dispositivo, red y concurrencia \cite{deloitte2024,adobe2026}.

\subsection{Abuso de concurrencia}

El servidor no mantiene sesiones activas ni contabiliza reproducciones simultáneas. Cada login genera un UUID independiente y no existe un límite por cuenta. Las solicitudes de licencia tampoco crean un registro de reproducción que pueda detenerse o expirar mediante \textit{heartbeat}. En consecuencia, una misma identidad puede iniciar varias reproducciones y obtener varias licencias sin recibir una respuesta distinta.

Esta carencia amplifica los efectos de la compartición de credenciales y del robo de tokens. Incluso si todas las licencias fueran criptográficamente válidas, la plataforma no podría determinar si corresponden al número de streams permitido por el plan. La defensa requiere un estado de sesión independiente del DRM, mantenido por la aplicación y consultado en los puntos de contenido y licencia.

\subsection{Ausencia de validaciones contextuales}

La decisión de Fase~0 se reduce a dos condiciones: que el token exista en memoria y que el activo figure en la lista de permisos. No se verifica que la petición conserve la misma IP, dispositivo o agente de usuario que el login. Tampoco se comprueba una sesión de reproducción activa, el ritmo de licencias, la relación entre KID y activo, el estado de riesgo o la concurrencia.

La ruta pública elimina incluso esas dos condiciones. El navegador externo pudo obtener dos respuestas de licencia y comenzar la reproducción sin enviar credenciales. Esta es la evidencia más directa de la hipótesis del TFM: el CDM y Widevine pueden funcionar correctamente y, al mismo tiempo, la plataforma puede autorizar de forma insuficiente el acceso al servicio.

La demostración de clave conocida añade un límite diferente. Al introducir manualmente el KID y la clave de laboratorio, Shaka Player reprodujo el activo CENC sin configurar servidor de licencias. Esta prueba no constituye un hallazgo contra Widevine, puesto que la clave no se extrajo del CDM. Ilustra que, una vez comprometido el material de contenido, las validaciones del servidor de licencias dejan de intervenir para ese activo.

\section{Escenarios de ataque reproducidos}

Los escenarios se diseñaron para responder a preguntas concretas y evitar pruebas ambiguas. Primero se verificó que el sistema pudiera distinguir un usuario permitido de otro sin \textit{entitlement}. Después se comprobó si esa decisión seguía siendo efectiva al abandonar el cliente oficial. Finalmente se evaluaron las rutas directas y el efecto de disponer de una clave conocida.

\subsection{Procedimiento}

El procedimiento aplicado fue el siguiente:

\begin{enumerate}
    \item Levantar los siete contenedores de Fase~0 y comprobar que cliente, CDN, origen y proxy respondían en sus puertos previstos.
    \item Iniciar sesión con el usuario autorizado, obtener el catálogo y reproducir el activo Widevine mediante \texttt{/platform/license}.
    \item Reproducir desde la misma plataforma el activo CENC local configurado con ClearKey.
    \item Iniciar sesión con el usuario sin permisos y comprobar el estado de ambos controles de reproducción.
    \item Abrir el reproductor externo sin iniciar sesión, mantener el MPD público y configurar \texttt{/license/no\_auth} como servidor Widevine.
    \item Repetir el escenario en el laboratorio DRM, registrar soporte EME, peticiones y tamaños de las respuestas de licencia.
    \item Solicitar el MPD local tanto a través de Varnish como directamente al origen publicado en \texttt{8081}.
    \item Introducir el KID y la clave de laboratorio en la demostración ClearKey y reproducir sin servidor de licencias.
    \item Solicitar al laboratorio que intentara mostrar claves Widevine y registrar la imposibilidad de exportarlas mediante EME.
\end{enumerate}

Cada resultado se contrastó con el estado del elemento de vídeo, la respuesta visible de la aplicación, las cabeceras HTTP y los logs de los contenedores. No se consideró exitoso un escenario basándose únicamente en mensajes generados por la propia interfaz.

\subsection{Resultados}

La Tabla~\ref{tab:resultados-fase0} recoge el resultado observado de cada escenario.

\begin{table}[H]
\centering
\caption{Resultados de los escenarios reproducidos en Fase~0}
\label{tab:resultados-fase0}
\begin{tabularx}{\textwidth}{p{0.09\textwidth}p{0.31\textwidth}Yp{0.14\textwidth}}
\toprule
\textbf{ID} & \textbf{Escenario} & \textbf{Resultado observado} & \textbf{Conclusión} \\
\midrule
E0-01 & Reproducción oficial autorizada & Widevine y ClearKey reproducen después de validar identidad y activo. & Control válido \\
E0-02 & Usuario sin \textit{entitlement} & Login correcto, catálogo denegado y botones deshabilitados. & Control válido \\
E0-03 & Reproductor Widevine externo & Reproducción iniciada con MPD y \texttt{no\_auth}, sin login. & Vulnerable \\
E0-04 & Acceso directo al origen & El MPD local responde por \texttt{8081} sin pasar por la CDN. & Vulnerable \\
E0-05 & Contenido local por CDN & MPD y segmentos accesibles sin token; la caché devuelve \texttt{HIT}. & Vulnerable \\
E0-06 & Reproducción con clave conocida & El activo CENC reproduce sin ninguna solicitud de licencia. & Impacto confirmado \\
E0-07 & Exportación de claves Widevine & El navegador no expone claves del CDM mediante EME. & No reproducido \\
\bottomrule
\end{tabularx}
\end{table}

El resultado E0-03 es el más relevante. El vídeo externo permaneció sin pausar, alcanzó estado de datos completos y avanzó sobre una duración aproximada de 888 segundos. El laboratorio repitió el comportamiento y registró dos respuestas de licencia. Por tanto, no se trató de una simulación visual, sino de una reproducción Widevine real fuera del cliente oficial.

E0-01 y E0-02 demuestran que la existencia de la vulnerabilidad no implica que toda la lógica de la plataforma sea inexistente. La autorización por catálogo funciona en el camino previsto. El problema es que puede eludirse seleccionando una ruta alternativa equivalente. Este resultado refuerza la necesidad de aplicar la misma política en todos los caminos capaces de entregar contenido o licencias.

\subsection{Evidencias}

Las evidencias principales obtenidas fueron las siguientes:

\begin{itemize}
    \item Estado \texttt{Reproduciendo: Sintel Widevine publico} en el cliente autorizado y avance real del elemento de vídeo.
    \item Dos botones deshabilitados y activos marcados como denegados para el usuario sin permisos.
    \item Evento \texttt{LICENSE\_REQUEST} dirigido a \texttt{/license/no\_auth} desde un cliente sin sesión.
    \item Respuestas binarias de licencia de 714 y 1229 bytes, seguidas de \texttt{PLAYBACK\_STARTED}.
    \item Elemento de vídeo externo sin pausar, con \texttt{readyState} completo y tiempo de reproducción creciente.
    \item Respuestas HTTP \texttt{200} para el MPD local tanto mediante \texttt{8080} como mediante \texttt{8081}.
    \item Cabecera \texttt{X-Cache: HIT} en la entrega del MPD a través de Varnish.
    \item Solicitudes registradas por Nginx para \texttt{video\_init.mp4} y segmentos \texttt{video\_N.m4s} durante la prueba ClearKey.
    \item Evento \texttt{PLAYBACK\_STARTED\_WITH\_KNOWN\_KEY} sin endpoint de licencias configurado.
    \item Evento \texttt{KEY\_EXTRACTION\_BLOCKED}, coherente con la imposibilidad de exportar claves Widevine mediante EME.
\end{itemize}

La evidencia también mostró aspectos secundarios de endurecimiento: uso de HTTP, CORS global, ausencia de cabeceras como CSP y publicación del token en el bloque de estado. Estos elementos no son necesarios para reproducir el bypass principal, pero amplían la superficie de exposición y deben corregirse en una arquitectura destinada a producción.

\section{Limitaciones del modelo DRM tradicional}

Los resultados confirman que Widevine cumple el cometido evaluado dentro del CDM: el navegador identifica el sistema de claves, genera un \textit{challenge}, procesa una licencia binaria y reproduce el contenido cifrado sin exponer la clave a JavaScript. La auditoría no halló una ruptura criptográfica de Widevine ni un método de extracción de claves. El fracaso se encuentra en la política que decide quién puede alcanzar el servicio de licencia y bajo qué contexto.

La ruta \texttt{/license/no\_auth} demuestra que una licencia técnicamente válida puede ser emitida dentro de una reproducción ilegítima para la plataforma. Del mismo modo, el acceso directo al origen muestra que el cifrado no evita el consumo de ancho de banda ni el uso de la infraestructura por terceros. Finalmente, la ausencia de sesiones y concurrencia permite que credenciales válidas se utilicen de una forma incompatible con la política de negocio sin que el DRM perciba ninguna anomalía.

La demostración ClearKey establece otro límite: ningún servidor de licencias puede revocar una clave que ya se encuentra fuera de su perímetro de control. En un sistema real, la robustez del CDM, la rotación de claves, las políticas de licencia y mecanismos complementarios como \textit{watermarking} ayudan a reducir el impacto, pero no convierten al cliente en una frontera absoluta. La literatura sobre Widevine L3 confirma además que el nivel de seguridad del dispositivo condiciona fuertemente la resistencia del extremo cliente \cite{patat2022,rafi2023mobiledrm}.

El laboratorio presenta limitaciones que deben conservarse al generalizar los resultados. Utiliza vídeo bajo demanda en lugar de ingesta en directo, un único navegador, usuarios de demostración, HTTP local y un servicio Widevine público. No reproduce carga masiva, diversidad geográfica, múltiples CDN ni políticas comerciales de licencia. Tampoco demuestra leeching completo de segmentos Widevine a través de la CDN local, puesto que el activo Widevine se aloja externamente. Estas limitaciones reducen la validez externa del experimento, pero no alteran el hallazgo central: si existe un camino de licencia sin autorización, la protección criptográfica puede seguir funcionando mientras el control de acceso de negocio queda anulado.

La auditoría permite derivar cinco requisitos para la arquitectura defensiva: eliminar rutas de licencia sin autenticación; impedir el acceso directo a origen y servidor de licencias; separar identidad y sesión de reproducción; emitir credenciales efímeras ligadas a usuario, activo y dispositivo; y registrar señales que permitan limitar concurrencia, detectar abuso y revocar acceso. El capítulo siguiente desarrolla estos requisitos como una arquitectura de defensa en profundidad, mientras que el capítulo de implementación describe su materialización progresiva en Fase~1 y Fase~2.
